\chapter*{Lời Mở Đầu}
\addcontentsline{toc}{chapter}{Lời Mở Đầu}
\markboth{Lời Mở Đầu}{Lời Mở Đầu}

\begin{truyen}{Tại sao là các cặp đối xứng?}{Human life is shaped by living between two poles.}
\chuhoa{C}{on người hiếm khi sống gọn trong một phía. Lúc thì muốn gần, lúc thì muốn tránh. Lúc muốn hơn người, lúc chỉ muốn yên thân. Lúc phải nói cho ra nhẽ, lúc lại phải ngậm lại kẻo vỡ thêm. Đời sống nó cứ giằng vậy, nên ai nhìn nó một màu thường sống rất mệt.}
\textit{(Human beings rarely live in one dimension. At times we want closeness, at times distance. We want to gain, yet fear losing. We need to speak to protect ourselves, but also need silence to listen and to repair ourselves. Life, therefore, is not only a story of choosing, but a story of balance.)}

Cuốn \textbf{100 Cặp Đối Xứng Triết Lý Nhân Sinh} được sắp thành bốn phần lớn. Phần một là những cặp giúp ta \textbf{nhìn đời}: \textbf{Có và Không}, \textbf{Được và Mất}, \textbf{Gần và Xa}, \textbf{Cho và Nhận}, \textbf{Giữ và Buông}. Phần hai là những cặp giúp ta \textbf{đi đường}: \textbf{Nhanh và Chậm}, \textbf{Mạnh và Yếu}, \textbf{Tiến và Lùi}, \textbf{Hơn và Kém}, \textbf{Thành và Bại}. Phần ba là những cặp giúp ta \textbf{xử thế}: \textbf{Đúng và Sai}, \textbf{Thật và Giả}, \textbf{Im và Nói}, \textbf{Nhớ và Quên}, \textbf{Cứng và Mềm}. Phần bốn là những cặp giúp ta \textbf{trưởng thành}: \textbf{Tin và Nghi}, \textbf{Mơ và Tỉnh}, \textbf{Mới và Cũ}, \textbf{Mở và Khép}, \textbf{Đi và Về}.
\textit{(\textbf{100 Symmetrical Pairs of Life Philosophy} is arranged into four main parts. The first offers pairs that help us \textbf{see life}: \textbf{Having and Not Having}, \textbf{Gain and Loss}, \textbf{Near and Far}, \textbf{Giving and Receiving}, \textbf{Holding and Letting Go}. The second offers pairs that help us \textbf{walk the road}: \textbf{Fast and Slow}, \textbf{Strong and Weak}, \textbf{Advance and Retreat}, \textbf{More and Less}, \textbf{Success and Failure}. The third offers pairs that help us \textbf{deal with others}: \textbf{Right and Wrong}, \textbf{Real and False}, \textbf{Silence and Speech}, \textbf{Remembering and Forgetting}, \textbf{Firmness and Softness}. The fourth offers pairs that help us \textbf{mature}: \textbf{Trust and Doubt}, \textbf{Dreaming and Awakening}, \textbf{New and Old}, \textbf{Opening and Closing}, \textbf{Going and Returning}.)}

Mỗi chương có năm tình huống nhỏ. Mỗi tình huống gồm một mẩu chuyện ngắn, một đoạn chốt thẳng, một câu tiếng Anh để nhớ nhanh, và vài câu hỏi để người đọc tự lật mình ra xem lại. Cách viết này giữ cho cuốn sách vừa đọc được như sách thường, vừa xài được trong lớp học, sinh hoạt, hay nói chuyện gia đình.
\textit{(Each chapter contains five small situations. Each situation includes a short story, a lesson, a quick English line to remember, and a few questions for reflection. This design allows the book to work both as a regular printed book and as material for classrooms, reading clubs, or family conversations.)}

Điều chính không phải là chốt xem bên nào thắng. Điều chính là thấy cho ra \textbf{mỗi bên đều có cái giá của nó}, để còn biết lúc nào nên nghiêng, lúc nào nên dừng, lúc nào nên rút chân lại trước khi quá muộn.
\textit{(The most important thing in this book is not deciding which side wins. More important is learning to see that \textbf{every pole carries its own price}, and to live in a way that does not let one side drag us too far.)}
\end{truyen}

\begin{baihoc}
Người trưởng thành không phải là người lúc nào cũng nói đạo lý chuẩn bài. Thường chỉ là người biết mình đang lệch ở đâu, và chịu chỉnh lại.
\textit{(A mature person is not one who always chooses one side, but one who knows when to lean, when to pause, and when to return to the center to keep balance.)}
\end{baihoc}

\begin{ghinhoanh}
\textbf{Short English to remember:}

Life is not flat. It moves between poles.

Balance is often wiser than victory.
\end{ghinhoanh}

\begin{tuhocnhanh}
1. Em thường nghiêng quá mức ở phía nào trong cách sống của mình? \textit{(Which side do you tend to lean too far toward in your way of living?)}

2. Nếu phải chọn một cặp để đọc trước, em sẽ chọn cặp nào vì nó đang đúng với hoàn cảnh của em? \textit{(If you had to choose one pair to read first, which would it be, and why does it match your current situation?)}
\end{tuhocnhanh}

\ngancach
